\chapter{Cross-model comparison}
\label{ch:crossmodel}

\begin{saybox}[title=What to say at the start of this section]
The point of this chapter is not to crown a winner. It is to show that the
original protocol could not rank models at all, and the revised one can.
That is a statement about the evaluation, and it is the reason the
evaluation had to change before any model-selection decision could be
trusted.
\end{saybox}

\section{The same models under the two protocols}

\begin{table}[htbp]\centering\small
\caption{Every evaluated configuration under both protocols. The
context-holdout column is the mean test $\rsq$ across the six specialists;
the leave-rule-out columns are the mean and median of the per-fold values,
averaged over the six specialists. Sorted by leave-rule-out median.}
\label{tab:both_protocols}
\begin{tabular}{lrrrr}
\toprule
Model & context-holdout $\rsq$ & LORO mean & LORO median & rank change \\
\midrule
LightGBM (production) & --- & 0.627 & \textbf{0.780} & --- \\
Random Forest (production) & --- & 0.582 & \textbf{0.791} & --- \\
XGBoost (production) & --- & 0.550 & \textbf{0.719} & --- \\
extra trees & 0.608 & 0.209 & 0.492 & $\downarrow$ \\
gradient boosting & 0.597 & 0.280 & 0.450 & $\downarrow$ \\
small MLP & 0.595 & 0.219 & 0.342 & $\downarrow$ \\
shallow random forest & 0.597 & $-0.230$ & 0.278 & $\downarrow\downarrow$ \\
shallow decision tree & 0.553 & $-1.733$ & 0.192 & $\downarrow\downarrow$ \\
RBF SVR & 0.557 & $-0.454$ & $-0.153$ & $\downarrow\downarrow$ \\
AdaBoost & 0.601 & $-1.248$ & $-0.178$ & $\downarrow\downarrow\downarrow$ \\
elastic net & 0.592 & $-0.772$ & $-0.675$ & $\downarrow\downarrow\downarrow$ \\
lasso & 0.595 & $-0.690$ & $-0.736$ & $\downarrow\downarrow\downarrow$ \\
Bayesian ridge & 0.597 & $-0.872$ & $-0.795$ & $\downarrow\downarrow\downarrow$ \\
ridge regression & 0.594 & $-0.874$ & $-0.798$ & $\downarrow\downarrow\downarrow$ \\
PLS regression & 0.594 & $-1.106$ & $-0.952$ & $\downarrow\downarrow\downarrow$ \\
$k$-NN ($k=5$) & 0.461 & $-3.412$ & $-1.537$ & $\downarrow\downarrow\downarrow$ \\
dummy (mean) & $-0.004$ & $-15.637$ & $-12.833$ & --- \\
\bottomrule
\end{tabular}
\end{table}

\begin{findingbox}[title=The discriminating power of the two protocols]
Under the original protocol, thirteen of the sixteen non-trivial
configurations fall within the band $[0.55, 0.61]$ --- a spread of six
hundredths of an $\rsq$ point separating a single shallow decision tree
from a 2000-round gradient ensemble.

Under leave-rule-out, the same configurations span from $+0.79$ to
$-1.54$ in median $\rsq$: a range more than twenty times wider. The
revised protocol does not merely give lower numbers; it gives
\emph{informative} numbers.
\end{findingbox}

\section{Per-specialist comparison for the production families}

\tbl{tab_loro_ensembles_by_specialist}{Leave-rule-out results by
specialist for the three production ensembles (repeated here for the
cross-model comparison).}

\begin{table}[htbp]\centering\small
\caption{Original context-holdout test $\rsq$ against pooled out-of-fold
$\rsq$ under leave-rule-out, per specialist and model. EBM appears in the
left column only; it has no leave-rule-out evaluation.}
\label{tab:protocol_delta}
\begin{tabular}{llrrr}
\toprule
Specialist & Model & context-holdout & pooled OOF (LORO) & $\Delta$ \\
\midrule
\multirow{3}{*}{soft / general}
 & Random Forest & 0.956 & 0.954 & $-0.002$ \\
 & LightGBM & 0.971 & 0.960 & $-0.011$ \\
 & XGBoost & 0.968 & 0.957 & $-0.011$ \\
\addlinespace[2pt]
\multirow{3}{*}{semi-hard / general}
 & Random Forest & 0.905 & 0.921 & $+0.016$ \\
 & LightGBM & 0.920 & 0.940 & $+0.020$ \\
 & XGBoost & 0.927 & 0.918 & $-0.009$ \\
\addlinespace[2pt]
\multirow{3}{*}{hard / general}
 & Random Forest & 0.927 & 0.582 & $\mathbf{-0.345}$ \\
 & LightGBM & 0.932 & 0.669 & $\mathbf{-0.263}$ \\
 & XGBoost & 0.933 & 0.665 & $\mathbf{-0.268}$ \\
\addlinespace[2pt]
\multirow{3}{*}{soft / safety}
 & Random Forest & 0.572 & 0.844 & $+0.272$ \\
 & LightGBM & 0.573 & 0.867 & $+0.294$ \\
 & XGBoost & 0.594 & 0.786 & $+0.192$ \\
\addlinespace[2pt]
\multirow{3}{*}{semi-hard / safety}
 & Random Forest & 0.342 & 0.755 & $+0.413$ \\
 & LightGBM & 0.303 & 0.811 & $+0.508$ \\
 & XGBoost & 0.317 & 0.655 & $+0.338$ \\
\addlinespace[2pt]
\multirow{3}{*}{hard / safety}
 & Random Forest & $-0.097$ & 0.642 & $+0.739$ \\
 & LightGBM & $-0.024$ & 0.730 & $+0.754$ \\
 & XGBoost & $-0.011$ & 0.593 & $+0.604$ \\
\bottomrule
\end{tabular}
\end{table}

\begin{findingbox}[title=The protocol change is not uniformly a penalty]
Only one specialist --- hard/general --- loses substantially, and the loss
is traced in \Cref{sec:fold-clostridia} to one held-out rule. The soft and
semi-hard quality specialists are essentially unchanged. All three safety
specialists \emph{gain}, by between $0.19$ and $0.75$ $\rsq$ points.

The honest summary is therefore not ``the real numbers are lower.'' It is
that a single random context split gave an unreliable reading in both
directions, and that a $56$-fold structured evaluation replaces it with a
distribution rather than a point.
\end{findingbox}

\section{Shared versus idiosyncratic failures}

\fig{fig34_model_agreement}{0.98}{%
  Pairwise agreement between the three ensembles at fold level
  (reproduced). $r = 0.877$--$0.939$.}

All three models rank the clostridia fold as their worst result, and all
three are positive on the majority of folds. The differences between them
are concentrated in the mid-range folds, where LightGBM's leaf-wise growth
appears to capture rule-to-rule differences that the depth-limited XGBoost
configuration does not.

\begin{limitbox}[title=A caution about attributing correlated errors]
Correlated failure across models is consistent with a shared data
limitation, but it is also consistent with a shared inductive bias: all
three are tree ensembles and all three inherit the extrapolation bound of
\cref{eq:rf-bound}. Distinguishing these would require a model family that
can extrapolate --- a parametric or mechanistic model --- evaluated on the
same folds. That experiment has not been run, and this report does not
claim the distinction has been made.
\end{limitbox}

\section{Recommended model under the revised protocol}

On the evidence assembled here, LightGBM is the appropriate production
choice: highest pooled out-of-fold $\rsq$ in all six specialists, highest
median fold $\rsq$ in five of six, and the smallest error on the single
hardest fold in the study. Random Forest is a close second with a
marginally higher median on two specialists and a more interpretable
uncertainty story.

That recommendation carries one explicit caveat: it is a recommendation
about performance on \emph{synthetic rule transfer}, which is the best
evidence currently available and is not the same thing as real-world
accuracy. \Cref{ch:external} presents what real-world evidence exists.
